从 10 个人里选 3 个当委员，再从这 3 人中指定一位主席，有多少种结果？可以先选委员会再挑主席，\(\binom{10}{3}\times3=360\)；也可以先定主席再补两名委员，\(10\times\binom92=360\)。两条路径的算式毫无相似之处，答案却分毫不差——因为它们数的是同一批东西。

这件小事透露了本章的立场。计数出错，几乎从不是乘法算错，而是``一个结果''没定义清楚：顺序算不算数，能不能重复，两个分支会不会撞在一起。定义一旦模糊，公式再熟也只是把错误算得更快。反过来，定义清楚之后，公式往往不必记——它可以由生成过程重新推出来。

真正贯穿全章的动作只有一个：把难数的集合对应到一个已经会数的集合上。对应可以是一一的，那就直接搬答案；可以是整齐的 \(d\) 对一，那就先数源集合再除以 \(d\)；也可以是换个方向重新清点同一批对象，两个结果必须相等，于是产生恒等式。当对应实在建立不起来时，还有两条退路：只断言碰撞必然存在而不给出位置，或者把重叠部分逐层加减干净。

\begin{center}
\begin{tikzpicture}[x=1cm,y=1cm,font=\small,>=stealth]
  \foreach \x/\a/\b in {0/{结果怎样}/{一步步生成}, 3.1/{换一种数法}/{两边必须相等}, 6.2/{只问会不会}/{撞在一起}, 9.3/{重叠部分}/{扣掉再补回}}{
    \draw[black!45,fill=black!4] (\x,-0.8) rectangle (\x+2.5,0.8);
    \node[align=center] at (\x+1.25,0) {\a\\\b};}
  \foreach \x in {2.5,5.6,8.7}{\draw[->,black!55] (\x,0) -- (\x+0.6,0);}
  \foreach \x/\s in {0/{第一节}, 3.1/{第二节}, 6.2/{第三节}, 9.3/{第四节}}{
    \node[black!55,font=\footnotesize] at (\x+1.25,-1.2) {\s};}
\end{tikzpicture}

\emph{四节各处理对应关系的一个侧面：怎样生成、怎样换算、对应失败时的必然碰撞、以及对应重叠时的修正。}
\end{center}

\hyperref[sec:04-Discrete-Mathematics/01-Combinatorics/01]{``加法、乘法与排列组合''}把加法、乘法、排列、组合收进同一条线索：互斥分支相加，整齐的连续选择相乘，而排列到组合、重复字母的排列都只是同一次除法——除掉一个可见结果背后的分身数。这一节也说清了乘法原理真正要求的是什么：每步的选项个数稳定，而不是选项本身相同。

\hyperref[sec:04-Discrete-Mathematics/01-Combinatorics/02]{``二项式与双重计数''}把``同一批对象换两种方式清点''升级成证明方法。二项式定理、Pascal 恒等式、对称性、Vandermonde 恒等式，各自对应一句组合叙述；代数展开只能核对，给不出理由。这一节还给数据工作留了一个直接可用的习惯：同一个指标换口径统计，两个数对不上，差额本身就是线索。

\hyperref[sec:04-Discrete-Mathematics/01-Combinatorics/03]{``鸽巢原理与必然碰撞''}处理另一类问题——不需要数出具体多少，只需要断言某种结构必定出现。原理的陈述近乎平凡，难度全部转移到设计上：谁是对象，什么是盒子，怎样让``同盒''恰好等于你想要的结论。哈希碰撞的必然性、通用无损压缩的不存在，都是这一下设计的产物。

\hyperref[sec:04-Discrete-Mathematics/01-Combinatorics/04]{``容斥原理与重复计数''}回到精确计数，面对的是分支重叠时该扣多少。交替符号不是口诀，它保证每个元素的净计数恰好为 1。同一条式子换成概率测度就是并事件公式，截断到第一层就是 union bound。

四节按依赖顺序排列，篇幅都不长，第一次读建议顺着走。若当下的困难很具体，也可以直接切入：分不清该加还是该乘，看第一节；想知道某个组合恒等式为什么成立，看第二节；只需要证明``某种冲突一定发生''，看第三节；统计口径重叠、去重结果对不上，看第四节。

读完之后可检验的变化是能做一次计数审计：说出``一个结果''由哪些信息完全确定，说出每个结果被生成了几遍，说出各分支是否互斥且穷尽。这三问答不上来，套用的公式就还没有可靠含义。往后学随机变量、离散分布与算法复杂度分析时，这套审计会被反复调用。
